% Generated draft: requires mathematical and source review.
\documentclass[11pt]{article}
\usepackage[margin=25mm]{geometry}
\usepackage{fontspec}
\setmainfont{Latin Modern Roman}
\usepackage{amsmath,amssymb,mathtools,unicode-math}
\setmathfont{Latin Modern Math}
\usepackage{enumitem,needspace,xurl,hyperref}
\setlength{\parindent}{0pt}
\setlength{\parskip}{5pt}
\setlength{\emergencystretch}{3em}
\newcommand{\N}{\mathbb{N}}
\newcommand{\Z}{\mathbb{Z}}
\newcommand{\Q}{\mathbb{Q}}
\newcommand{\R}{\mathbb{R}}
\newcommand{\C}{\mathbb{C}}
\newcommand{\F}{\mathbb{F}}
\newcommand{\A}{\mathbb{A}}
\newcommand{\PP}{\mathbb{P}}
\newcommand{\E}{\mathbb{E}}
\newcommand{\eps}{\varepsilon}
\newcommand{\dd}{\,\mathrm{d}}
\begin{document}
\section*{UnsolvedMath ID: 800697}
\textbf{Source problem code:} catalogue entry
\subsection{Definitions and mathematical statement}
Let \(F\) be a field. Its Brauer group \(\operatorname{Br}(F)\) consists of Brauer equivalence classes of finite-dimensional central simple \(F\)-algebras, with addition induced by tensor product. For a positive integer \(n\), write
\[
\operatorname{Br}(F)[n]=\{\alpha\in\operatorname{Br}(F):n\alpha=0\}.
\]
A cyclic algebra of degree \(d\) over \(F\) is an algebra \((L/F,\sigma,a)\), where \(L/F\) is a cyclic Galois extension of degree \(d\), \(\sigma\) generates its Galois group, and \(a\in F^\times\). It is generated by \(L\) and an element \(u\), subject to \(u^d=a\) and \(u\ell=\sigma(\ell)u\) for all \(\ell\in L\).

The source asks whether its subgroup \(\operatorname{Br}_p[n](F)\) is generated by classes of cyclic algebras whose degrees divide \(n\). With the source's notation retained, the requested assertion is
\[
\operatorname{Br}_p[n](F)
\stackrel{?}{=}
\left\langle [A]\in\operatorname{Br}(F):
A\text{ is a cyclic }F\text{-algebra and }\deg(A)\mid n
\right\rangle.
\]
Here generation means taking finite sums and additive inverses in the Brauer group. The supplied record does not define the subscript \(p\), specify its relation to \(n\) or \(\operatorname{char}(F)\), or state the intended restrictions on \(F\) and \(n\). Consequently, the displayed equality records the question but does not supply a fully specified universally quantified conjecture. In particular, identifying \(\operatorname{Br}_p[n](F)\) with \(\operatorname{Br}(F)[n]\) would require information absent from the record.
\subsection{Short English statement}
Can every class in the source's \(n\)-torsion subgroup \(\operatorname{Br}_p[n](F)\) be generated by classes of cyclic algebras of degree dividing \(n\)? The meaning of \(p\) and the intended field and integer hypotheses remain unspecified.
\subsection{Sources}
\begin{itemize}
\item Local UnsolvedMath registry: permanent ID \(800697\), source problem code \texttt{catalogue entry}, titled ``Open Problems on Central Simple Algebras --- Problem 7.''
\item Auel, Brussel, Garibaldi, and Vishne, \emph{Open Problems on Central Simple Algebras} (2010), Problem 7. Supplied source URL: \url{https://arxiv.org/abs/1006.3304}. The record reports access on 2026-10-16 and extraction from arXiv LaTeX; these are supplied provenance details, not independently verified here.
\item The statement includes the attribution \texttt{[Albert:cyclics, p. 126]}. The supplied record does not resolve this bibliographic key to a complete reference.
\end{itemize}
\end{document}
